\documentclass[11pt,a4paper]{article}

% --- PAQUETES DE CONFIGURACIÓN ---
\usepackage[spanish]{babel}
\usepackage[utf8]{inputenc}
\usepackage[T1]{fontenc}
\usepackage{geometry}
\geometry{margin=2.5cm}
\usepackage{hyperref}
\usepackage{xcolor}
\usepackage{listings}
\usepackage{booktabs}
\usepackage{titlesec}

% --- CONFIGURACIÓN DE COLORES PARA CÓDIGO ---
\definecolor{codegreen}{rgb}{0,0.6,0}
\definecolor{codegray}{rgb}{0.5,0.5,0.5}
\definecolor{codepurple}{rgb}{0.58,0,0.82}
\definecolor{backcolour}{rgb}{0.96,0.96,0.96}

% --- ESTILO DE LISTINGS (PYTHON) ---
\lstdefinestyle{pystyle}{
    backgroundcolor=\color{backcolour},   
    commentstyle=\color{codegreen},
    keywordstyle=\color{magenta},
    numberstyle=\tiny\color{codegray},
    stringstyle=\color{codepurple},
    basicstyle=\ttfamily\small,
    breakatwhitespace=false,         
    breaklines=true,                 
    captionpos=b,                    
    keepspaces=true,                 
    numbers=left,                    
    numbersep=5pt,                  
    showspaces=false,                
    showstringspaces=false,
    showtabs=false,                  
    tabsize=4
}
\lstset{style=pystyle, extendedchars=true, literate={á}{{\'a}}1 {é}{{\'e}}1 {í}{{\'i}}1 {ó}{{\'o}}1 {ú}{{\'u}}1 {ñ}{{\~n}}1 {Ó}{{\'O}}1
}

% --- DATOS DEL DOCUMENTO ---
\title{\textbf{Manual de Referencia: Profiling y Optimización Multihilo de Simuladores Físicos en Python}}
\author{Guía Técnica de Desarrollo}
\date{\today}

\begin{document}

\maketitle
\tableofcontents
\newpage

% ==============================================================================
\section{Evaluación y Profiling del Tiempo de Ejecución}
% ==============================================================================
Cuando un simulador físico escrito en Python incluye un cálculo numérico (resolución de ODEs), una interfaz gráfica (GUI) y almacenamiento en disco (I/O), es fundamental medir qué porcentaje del tiempo total consume cada componente para detectar cuellos de botella.

\subsection{Medición Directa por Bloques con \texttt{time.perf\_counter()}}
La forma más precisa para aislar estos bloques dentro del bucle de integración principal consiste en medir el tiempo acumulado con un temporizador de alta resolución del sistema operativo.

\begin{lstlisting}[language=Python, caption={Medición del tiempo por bloques en el bucle principal.}]
import time

# Acumuladores de tiempo total (en segundos)
t_sim, t_gui, t_io = 0.0, 0.0, 0.0
n_pasos = 0

def paso_de_simulacion():
    global t_sim, t_gui, t_io, n_pasos
    
    # 1. Cálculo Numérico
    t0 = time.perf_counter()
    estado_nuevo = resolver_ode_paso(estado_actual)
    t1 = time.perf_counter()
    t_sim += (t1 - t0)
    
    # 2. Escritura en Disco
    t0 = time.perf_counter()
    escribir_archivo(estado_nuevo)
    t1 = time.perf_counter()
    t_io += (t1 - t0)
    
    # 3. Actualización de Interfaz / Gráfico
    t0 = time.perf_counter()
    actualizar_grafico_gui(estado_nuevo)
    t1 = time.perf_counter()
    t_gui += (t1 - t0)
    
    n_pasos += 1

def mostrar_reporte():
    t_total = t_sim + t_gui + t_io
    print(f"--- REPORTE DE RENDIMIENTO ({n_pasos} pasos) ---")
    print(f"Simulación (Cálculo): {t_sim:.4f} s ({t_sim/t_total*100:.1f}%)")
    print(f"Escritura en Disco:   {t_io:.4f} s ({t_io/t_total*100:.1f}%)")
    print(f"Interfaz Gráfica:     {t_gui:.4f} s ({t_gui/t_total*100:.1f}%)")
    print(f"Tiempo Total Medido:  {t_total:.4f} s")
\end{lstlisting}

\subsection{Profiling Global con \texttt{cProfile} y Visualización}
Para inspeccionar de manera interna qué llamadas dentro de bibliotecas como Matplotlib, Scipy o la GUI consumen recursos, se ejecuta:

\begin{lstlisting}[language=bash, caption={Uso de cProfile y SnakeViz desde la terminal.}]
# Generar el archivo de profiling
python -m cProfile -o resultado.prof tu_programa.py

# Visualizar interactivamente en el navegador
pip install snakeviz
snakeviz resultado.prof
\end{lstlisting}

% ==============================================================================
\section{Estrategias para Desacoplar Renderizado e I/O}
% ==============================================================================
El principio clave para evitar que la GUI y la escritura en disco frenen la simulación es \textbf{desacoplar los tiempos de cada proceso}.

\subsection{Arquitectura Multihilo y Colas Thread-Safe}
\begin{itemize}
    \item \textbf{Hilo de Simulación:} Resuelve las ecuaciones diferenciales en un bucle continuo sin pausas.
    \item \textbf{Hilo Principal (GUI):} Mantiene la respuesta de la interfaz a ~60 FPS consumiendo solo estados recientes.
    \item \textbf{Hilo de Disco (I/O):} Recibe el 100\% de las muestras y las escribe en bloques.
    \item \textbf{Colas (\texttt{queue.Queue}):} Permiten la comunicación no bloqueante (\texttt{put\_nowait()}).
\end{itemize}

\subsection{Optimización de Disco y Gráficos}
\begin{enumerate}
    \item \textbf{Escritura en Bloques (\textit{Batching}):} Acumular datos en memoria y volcarlos a disco cada $N$ iteraciones mediante \texttt{writerows()} para minimizar las llamadas al sistema.
    \item \textbf{Subsampling (Decimación):} Reducir la tasa de actualización de la GUI (ej. enviar a la cola 1 de cada 10 pasos).
    \item \textbf{Blitting en Matplotlib:} Redibujar únicamente la línea modificada en lugar de recalcular ejes, grillas y etiquetas en cada cuadro.
\end{enumerate}

\begin{table}[h]
\centering
\caption{Resumen de estrategias de optimización.}
\begin{tabular}{lll}
\toprule
\textbf{Técnica} & \textbf{Problema que Resuelve} & \textbf{Impacto Esperado} \\
\midrule
Multihilo + \texttt{Queue} & Congelamiento de la simulación & Alto (solver independiente) \\
\textit{Batching} en I/O & Latencia por llamadas a sistema & Reduce el impacto de I/O a $<$1\% \\
Blitting / PyQtGraph & Sobrecarga de renderizado de la GPU/CPU & Muy alto (fluidez visual a 60 FPS) \\
\bottomrule
\end{tabular}
\end{table}

\newpage
% ==============================================================================
\section{Implementación Completa: Ejemplo Integrado}
% ==============================================================================
El siguiente código implementa un simulador completo en \textbf{PyQt6/PyQt5} utilizando tres hilos de ejecución concurrentes:
\begin{enumerate}
    \item \texttt{SimulationWorker} (Cálculo numérico).
    \item \texttt{DiskWorker} (Escritura en disco con \textit{batching}).
    \item \texttt{MainWindow} (Hilo principal de la GUI con \textit{Blitting} en Matplotlib).
\end{enumerate}

\begin{lstlisting}[language=Python, caption={Código completo del simulador desacoplado multihilo.}]
import sys
import time
import csv
import numpy as np
from queue import Queue, Empty

from PyQt6.QtCore import QThread, QTimer
from PyQt6.QtWidgets import (QApplication, QMainWindow, QWidget, 
                             QVBoxLayout, QHBoxLayout, QPushButton)

from matplotlib.backends.backend_qtagg import FigureCanvasQTAgg as FigureCanvas
from matplotlib.figure import Figure

# ==============================================================================
# 1. HILO DE ESCRITURA EN DISCO (I/O Worker con Batching)
# ==============================================================================
class DiskWorker(QThread):
    def __init__(self, disk_queue: Queue, filename="simulacion_resultados.csv", batch_size=1000):
        super().__init__()
        self.disk_queue = disk_queue
        self.filename = filename
        self.batch_size = batch_size
        self.running = False

    def run(self):
        self.running = True
        buffer = []

        with open(self.filename, mode="w", newline="", encoding="utf-8") as f:
            writer = csv.writer(f)
            writer.writerow(["Tiempo_s", "Posicion_x", "Velocidad_v"])

            while self.running or not self.disk_queue.empty():
                try:
                    data = self.disk_queue.get(timeout=0.1)
                    buffer.append(data)
                    self.disk_queue.task_done()

                    # Volcado masivo a disco al acumular N elementos
                    if len(buffer) >= self.batch_size:
                        writer.writerows(buffer)
                        buffer.clear()

                except Empty:
                    continue

            if buffer:
                writer.writerows(buffer)
                buffer.clear()

    def stop(self):
        self.running = False
        self.wait()

# ==============================================================================
# 2. HILO DE SIMULACIÓN (Cálculo Numérico Puro)
# ==============================================================================
class SimulationWorker(QThread):
    def __init__(self, gui_queue: Queue, disk_queue: Queue):
        super().__init__()
        self.gui_queue = gui_queue
        self.disk_queue = disk_queue
        self.running = False

    def run(self):
        self.running = True
        
        t = 0.0
        dt = 0.001  
        x = 1.0     
        v = 0.0     
        gamma = 0.1
        omega0 = 2.0 * np.pi

        step_counter = 0

        while self.running:
            # Integración numérica (Euler-Cromer)
            a = -2 * gamma * v - (omega0**2) * x
            v += a * dt
            x += v * dt
            t += dt

            step_counter += 1

            # 1. ENVIAR A DISCO: Se guardan TODOS los pasos
            try:
                self.disk_queue.put_nowait((t, x, v))
            except Exception:
                pass

            # 2. ENVIAR A GUI: Subsampling (1 de cada 10 pasos)
            if step_counter % 10 == 0:
                try:
                    self.gui_queue.put_nowait((t, x))
                except Exception:
                    pass

            time.sleep(0.0001)

    def stop(self):
        self.running = False
        self.wait()

# ==============================================================================
# 3. VENTANA PRINCIPAL (GUI y Renderizado con Blitting)
# ==============================================================================
class MainWindow(QMainWindow):
    def __init__(self):
        super().__init__()
        self.setWindowTitle("Simulador Multihilo (Cálculo + GUI + Disco)")
        self.resize(800, 500)

        self.gui_queue = Queue(maxsize=100)
        self.disk_queue = Queue(maxsize=50000)

        self.sim_worker = SimulationWorker(self.gui_queue, self.disk_queue)
        self.disk_worker = DiskWorker(self.disk_queue)

        main_widget = QWidget()
        self.setCentralWidget(main_widget)
        layout = QHBoxLayout(main_widget)

        controls_layout = QVBoxLayout()
        self.btn_start = QPushButton("Iniciar Simulación")
        self.btn_start.clicked.connect(self.toggle_simulation)
        controls_layout.addWidget(self.btn_start)
        controls_layout.addStretch()
        layout.addLayout(controls_layout)

        self.fig = Figure(figsize=(6, 4))
        self.canvas = FigureCanvas(self.fig)
        layout.addWidget(self.canvas)

        self.ax = self.fig.add_subplot(111)
        self.ax.set_xlim(0, 10)
        self.ax.set_ylim(-1.5, 1.5)
        self.ax.set_xlabel("Tiempo (s)")
        self.ax.set_ylabel("Posición (x)")
        self.ax.grid(True)

        (self.line,) = self.ax.plot([], [], 'b-', animated=True)

        self.t_data = []
        self.x_data = []
        self.background = None

        self.canvas.mpl_connect("draw_event", self.on_draw)

        self.timer = QTimer()
        self.timer.setInterval(16)  # ~60 FPS
        self.timer.timeout.connect(self.update_gui)

    def on_draw(self, event):
        self.background = self.canvas.copy_from_bbox(self.ax.bbox)
        self.ax.draw_artist(self.line)

    def toggle_simulation(self):
        if not self.sim_worker.isRunning():
            self.t_data.clear()
            self.x_data.clear()
            
            self.disk_worker.start()
            self.sim_worker.start()
            self.timer.start()
            
            self.btn_start.setText("Detener Simulación")
        else:
            self.sim_worker.stop()
            self.disk_worker.stop()
            self.timer.stop()
            
            self.btn_start.setText("Iniciar Simulación")

    def update_gui(self):
        if self.background is None:
            return

        updated = False
        while True:
            try:
                t, x = self.gui_queue.get_nowait()
                self.t_data.append(t)
                self.x_data.append(x)
                updated = True
            except Empty:
                break

        if updated:
            if self.t_data[-1] > 10:
                self.ax.set_xlim(self.t_data[-1] - 10, self.t_data[-1])
                self.canvas.draw()
                return

            # Renderizado Blitting
            self.line.set_data(self.t_data, self.x_data)
            self.canvas.restore_region(self.background)
            self.ax.draw_artist(self.line)
            self.canvas.blit(self.ax.bbox)

    def closeEvent(self, event):
        self.sim_worker.stop()
        self.disk_worker.stop()
        event.accept()

# ==============================================================================
# 4. PUNTO DE ENTRADA
# ==============================================================================
if __name__ == "__main__":
    app = QApplication(sys.argv)
    window = MainWindow()
    window.show()
    sys.exit(app.exec())
\end{lstlisting}

\end{document}
